\section{Discussion}
\label{sec:discuss}
% We introduce SWE-agent, an agent composed of an LM and ACI capable of autonomously solving software engineering tasks.
% SWE-agent achieves a state-of-the-art \num{12.5}\% performance on SWE-bench and \num{87.7}\% pass$@$\num{1} rate on HumanEvalFix.
% Through our design methodology, empirical results, and analysis, we demonstrate the process and value of designing agent-computer interfaces (ACIs) tailored to leverage LMs' strengths and mitigate their weaknesses.
% We release our code, prompts and generations, and feedback and agent history formats.
% The SWE-agent ACI is written to be extensible, and we look forward to the investigations on ACI design for software engineering that our contribution aims to motivate.
% We hope that SWE-agent will serve as a foundation that inspires future work towards more versatile and powerful agents.

We introduce SWE-agent, an agent composed of an LM and ACI capable of autonomously solving software engineering tasks.
% SWE-agent achieves a state-of-the-art \num{12.5}\% performance on SWE-bench and \num{87.7}\% pass$@$\num{1} rate on HumanEvalFix.
Through our design methodology, results, and analysis, we demonstrate the value of ACIs tailored to leverage LMs' strengths and mitigate their weaknesses.
Beyond empirical applications, we hope the further study of ACIs can also make principled use of and contribute to our understanding of language models and agents, analogous to the synergy between human-computer interaction (HCI) and psychology~\citep{carroll1997human}.
Humans and LMs have different characteristics, training objectives, specialities, and limitations~\citep{griffiths2020understanding,mccoy2023embers}, and the interaction design processes can be seen as systematic behavioral experimentation that could reveal more insights into these differences towards establishing a comparative understanding of human and artificial intelligence.
